\documentclass[12pt]{article}
\usepackage[T2A]{fontenc}
\usepackage[utf8]{inputenc}
\usepackage[russian]{babel}
\usepackage{amsmath, amssymb, amsfonts}
\usepackage{geometry}
\usepackage{fancyhdr}
\usepackage{microtype}
\usepackage{titlesec} % Пакет для управления стилем заголовков
\usepackage{enumitem}
\usepackage{tikz}
\usepackage{tikz-3dplot}
\usetikzlibrary{calc, intersections, through, backgrounds, arrows.meta}

% Настройка Section: по центру, жирный, крупный
\pagestyle{fancy}
% Настройка полей
\geometry{a5paper, margin=1cm, top=2cm, headheight=15pt}

% --- ЖЕСТКАЯ КОРРЕКЦИЯ ЦЕНТРИРОВАНИЯ ---
\pagestyle{fancy}
\fancyhf{} % Полная очистка всех полей (убирает невидимые отступы)
\fancyhead[C]{\thepage} % Номер строго по центру
\renewcommand{\headrulewidth}{0pt} % Убираем линию
\fancyfoot{} % Низ пустой

% Убираем внутренние отступы колонтитулов, которые могут давать перекос
\setlength{\headwidth}{\textwidth} 

% Переопределяем системный стиль plain, чтобы он не сбивал настройки
\makeatletter
\let\ps@plain\ps@fancy
\makeatother

\parindent=0pt
\parskip=0.6em
\titleformat{\section}
  {\normalfont\bfseries\centering}{}{0pt}{}

% Настройка Subsection: отступ 1см, жирный, обычный размер
% {команда}[формат]{стиль}{метка}{отступ_от_метки}{перед_заголовком}[после_заголовка]
\titleformat{\subsection}
  {\normalfont\bfseries}{\hspace{1cm}}{1pt}{}
\titlespacing*{\subsection}{1cm}{2ex plus 1ex minus .2ex}{0.5ex plus .2ex}

% Убираем автоматическую нумерацию, чтобы заголовки были как в оригинале
\setcounter{secnumdepth}{0}
\AddEnumerateCounter{\asbuk}{\russian@asbuk@alph}{щ} % section 8.1
\setlist{wide, nosep, font=\rmfamily\normalfont}
\setlist[enumerate, 2]{label=\asbuk*)}
\setlist[enumerate, 3]{label=\asbuk*)}


\begin{document}

\section{Глава 4. Квадрики}

\subsection{4.1. Кривые второго порядка}

\subsubsection*{4.1.1. Общие свойства кривых второго порядка (коники)}

\begin{enumerate}

	\item Рассмотрим семейство коник, зависящее от параметра $\lambda$:
	\[
		\frac{x^2}{a - \lambda} + \frac{y^2}{b - \lambda} = 1.
	\]
	Доказать, что:
	\begin{enumerate}
		\item фокусы всех этих коник совпадают;
		\item через каждую точку плоскости, не лежащую на осях координат, проходит ровно две коники данного семейства;
		\item в точках пересечения коники ортогнальны.
	\end{enumerate}

	\item Рассмотрим все правильные треугольники, вписанные в данную конику. Доказать, что их центры лежат на некоторой конике.

\end{enumerate}
\begin{enumerate}[resume*]
\subsubsection*{4.1.2. Теорема о пучке коник, проходящих через 4 точки}
	\item Предположим, что никакие 3 из точек $A$, $B$, $C$ и $D$ не лежат на одной прямой.
	\begin{enumerate}
		\item Пусть $f = 0$ -- уравнение некоторой коники, проходящей через 4 точки $A$, $B$, $C$ и $D$. Доказать, что
		\[
			f = \lambda AB \cdot CD + \mu AD \cdot BC,
		\]
		где $AB = 0$ -- уравнение, задающее прямую $AB$(точнее говоря, одно из этих уравненй).

		\item Пусть $f = 0$ и $g = 0$ -- уравнения двух различных коник, проходящих через точки $A$, $B$, $C$ и $D$. Доказать, что если $h = 0$ -- уравнение любой коники, проходящей через эти точки, то $h = \lambda f + \mu g$.
	\end{enumerate}

	\item Две коники имеют 4 общих точки. Доказать, что на эти точки лежат одной окружности тогда и только тогда, когда оси коник перпендикулярны.

	\item
	\begin{enumerate}
		\item Доказать, что кривая $ax^2 + bxy + cy^2 + \ldots = 0$ является гиперболой с перпендикулярными асимптотами тогда и только тогда, когда $a + c = 0$.
		\item Пусть $H$ -- точка пересечения высот треугольника $ABC$. Доказать, что любая коника, проходящая через точки $A$, $B$, $C$ и $H$, является гиперболой с перпендикулярными асимптотами.
	\end{enumerate}

	\item
	\begin{enumerate}
		\item (Теорема о бабочке.) Хорды $KL$ и $MN$ проходят через середину $O$ хорды $AB$ окружности. Доказать, что прямые $KN$ и $ML$ пересекают прямую $AB$ в точках, равноудаленных от точки $O$.
		\item Три коники имеют 4 общих точки. Прямая $AB$ высекает на двух из них хорды, имеющие общую середину $O$. Доказать, что $O$ является также серединой хорды, высекаемой прямой $AB$ на третьей конике.
	\end{enumerate}

	\item Для решения уравнения $4$й степени $x^4 + ax^3 + bx^2 + cx + d = 0$ достаточно найти точки пересечения кривых
	\[
		f = y - x^2 = 0 \quad \text{и} \quad g = y^2 + axy + by + cx + d = 0.
	\]
	Вторую кривую можно заменить на $\lambda f + g = 0$. Если эта кривая вырожденная, т.е. уравнение распадается на два линейных уравнения, то дальнейшее решение очевидно. Доказать, что для нахождения $\lambda$ нужно решить кубическое уравнение.

	\item
	\begin{enumerate}
		\item Шестиугольник $ABCDEF$ вписан в конику. Доказать, что точки пересчения прямых $AB$ и $DE$, $BC$ и $EF$, $CD$ и $FA$ лежат на одной прямой (Паскаль). 
		\item Доказать, что прямые Паскаля шестиугольников $ABCDEF$,

			$ADEBDC$ и $ADCFEC$ пересекаются в одной точке (Штейнер).
		\item Доказать, что прямые Паскаля шестиугольников $ABFDCE$,

			$AEFBDC$ и $ABDFEC$ пересекаются в одной точке (Кирхман).
		\item Доказать, что на прямой Паскаля лежит одна точка Штейнера и три точки Кирхмана.
	\end{enumerate}

\end{enumerate}

\subsubsection*{4.1.3. Парабола}

\begin{enumerate}

	\item Доказать, что касательная к параболе $y = x^2$ в точке $(x_0, y_0)$ пересекает ось $Oy$ в точке $(0, -y_0)$.

	\item Доказать, что середины параллельных хорд параболы $y = x^2$ лежат на одной прямой, параллельной оси $Oy$.

	\item На плоскости нарисована парабола $y = x^2$, но оси координат не нарисованы. С помощью циркуля и линейки построить оси координат.

	\item
	\begin{enumerate}
		\item Доказать, что точки параболы $x^2 = 4ay$ ранвоудалены от точки $(0, a)$ и прямой $y = -a$. Эти точка и прямая называются фокусом и директрисой параболы.
		\item Доказать, что фокус и директриса параболы опредлеены однозначно.
	\end{enumerate}

	\item Пусть $F$ -- фокус параболы, $d$ -- директриса. Доказать, что касательная в точке $A$ параболы является биссектрисой угла $FAH$, где $AH$ -- перпендикуляр к прямой $d$.

	\textit{Замечание.} Эта задача допускает следуюущую переформулировку: "Пучок лучей света, параллельных оси параболы, после отражения от параболы сходится в фокусе."

	\item Доказать, что парабола о фокусом $(a, 0)$ и директрисой $x = -a$ задается в полярных координатах уравнением $r = \frac{2a}{1 - \cos\varphi}.$

	\item Доказать, что преобразование $z \mapsto 2(1-z)^3$ переводит единичную окружность $|z| = 1$ в параболу.
	
	\item Существует ли преобразование плоскости, переводящее семейство парабол $x = a _ by _ y^2$ в семейство прямых?

	\item Окружность пересекает параболу в четырех точках. Доказать, что центр масс этих точек лежит на оси параболы.

	\item Две параболы, оси которых перпендикулярны, пересекаются в четырех точках. Доказать, что эти точки лежат на одной окружности.

	\item Хорда $P_1 P_2$ параболы проходит через фокус $F$. Доказать, что величина $P_1F^{-1} + P_2F^{-1}$ не зависит от выбора хорды.

	\item Из точки $O$ проведены касательные $OA$ и $OB$ к параболе с фокусом $F$. Доказать, что $\angle AFB = 2\angle AOB$, причем луч $OF$ -- биссектриса угла $AOB$.

	\item Доказать, что касательные $OA$ и $OB$ к параболе перпендикулярны тогда и только тогда, когда выполнено одно из следующих эквивалентных условий:
	\begin{enumerate}
		\item отрезок $AB$ проходит через фокус параболы;
		\item точка $O$ лежит на директрисе.
	\end{enumerate}

	\item Доказать, что касательные к параболе $x^2 = 4y$ в точках $(2t_i, t_i^2)$, $i = 1, 2$, пересекаются в точках $(t_1 + t_2, t_1 t_2)$.

	\item Касательные к параболе в точках $\alpha, \beta, \gamma$ образуют треугольник $ABC$. Доказать, что:
	\begin{enumerate}
		\item описанная окружность треугольника $ABC$ проходит через фокус параболы;
		\item высоты треугольника $ABC$ пересекаются в точке, лежащей на директрисе.
		\item $S_{ABC} = \frac{1}{2} S_{\alpha \beta \gamma}$
	\end{enumerate}

	\item Прямая $l$ получена из директрисы гомотеттие с центром в фокусе параболы и коэффициентом 2. Из точки $O \in l$ проведены касательные $OA$ и $OB$ к параболе. Доказать, что ортоцентром треугольника $AOB$ служит вершина параболы.

	% \item \textit{Указание.} Касательная. Прямая $y = ax + b$ пересекает параболу $y = x^2$ в точках $(x_1, y_1)$ и $(x_2, y_2)$, где $ax_i + b = x_i^2$. \textit{Указание.} Пусть прямая $l$ проходит через серединам двух параллельных хорд параболы. Ось $Oy$ перпендикулярна к параболе, пересекает ось $Oy$ в точке $B(0, -4ax_0^2)$.

	% Расстояние от точки $(x_0, 4ax_0^2)$ до директрисы равно $y_0 + a = 4ax_0^2 + a$.
	%
	% Пусть $v_0$ таков, что $b = 1, a = 0, x_0 = 0$. Тогда $AH = 4a_0x_0^2 + a = (px + q)^2 + r = a$.
	%
	% \textit{Указание.} Расстояние от точки $(r, \varphi)$ до директрисы равно $r \cos\varphi + 2a$, а параллелограмм от точки $B(0, -4ax_0^2)$ до директрисы равно $r \cos\varphi + 2a$, поэтому
	% \[
	% 	r = r\cos\varphi + 2a.
	% \]
	%
\end{enumerate}

\subsubsection*{4.1.4. Эллипс}

\begin{enumerate}

	\item Доказать, что для эллипса $\frac{x^2}{a^2} + \frac{y^2}{b^2} = 1$ существуют такие точка $F$ и прямая $d$, что отношение расстояний от точки эллипса до точки $F$ к расстоянию до прямой $d$ равно постоянному числу $e < 1$. Найти координаты точки $F$ и уравнение прямой $d$.

	\item Доказать, что сумма расстояний от точки $X$ до двух фокусов $F_1$ и $F_2$ постоянна.

	\item Доказать, что в полярных координатах с началом в фокусе $F_1$ и осью координат $F_1 F_2$ эллипс имеет уравнение $r = \frac{p}{1 - e\cos\varphi}$, где $e < 1$.

	\item
	\begin{enumerate}
		\item Найти уравнение касательной к эллипсу $\frac{x^2}{a^2} + \frac{y^2}{b^2} = 1$ в точке $(x_0, y_0)$.
		\item Прямая касается эллипса в точке $A$. Доказать, что она образует равные углые с отрезками $F_1A$ и $F_2A$.
	\end{enumerate}

	\item Доказать, что эллипс $\frac{x^2}{a^2} + \frac{y^2}{b^2} = 1$ является образом окружности $x^2 + y^2 = 1$ при линейном преобразовании.

	\item
	\begin{enumerate}
		\item Доказать, что все середины параллельных хорд эллипса лежат на одной хорде, проходящей через центр эллипса.
		\item Пусть хорда $AB$ эллипса проходит через его центр и середины всех хорд, параллельных $AB$, лежат на хорде $CD$. Доказать, что середины хорд, параллельных $CD$ лежат на $AB$ (Такие хорды $AB$ и $CD$ называются сопряженными диаметрами эллипса).
	\end{enumerate}

	\item
	\begin{enumerate}
		\item Доказать, что для любого параллелограмма существует эллипс, касающийся сторон параллелограмма в их серединах.
		\item Доказать, что для любого треугольника существует эллипс, касающийся сторон треугольника в их серединах.
	\end{enumerate}

	\item Пусть $AA'$ и $BB'$ -- сопряженные диаметры эллипса с центром $O$. Доказать, что:
	\begin{enumerate}
		\item площадь треугольника $AOB$ не зависит от выбора сопряженных диаметров;
		\item величина $OA^2 + OB^2$ не зависит от выбора сопряженных диаметров.
	\end{enumerate}

	\item
	\begin{enumerate}
		\item Доказать, что проекции фокусов эллипса на все касательные лежат на одной окружности.
		\item Пусть $d_1$ и $d_2$ -- расстояния от фокусов эллипса до касательной. Доказать, что величина $d_1 d_2$ не зависит от выбора касательной.
	\end{enumerate}

	\item Из точки $O$ проведены касательные $OA$ и $OB$ к эллипсу с фокусами $F_1$ и $F_2$. Доказать, что $\angle AOF_1 = \angle BOF_2$, $\angle AF_1O = \angle BF_1O$.

	\item Вокруг эллипса описан прямоугольник. Доказать, что длина его диагонали не зависят от положения прямоугольника.

	\item Нормаль к эллипсу в точке $A$ пересекает малую полуось в точке $Q$, $P$ -- проекция центра эллипса на нормаль. Доказать, что $AP \cdot AQ = a^2$, где $a$ -- большая полуось.

	\item Пусть $AA'$ и $BB'$ -- сопряженные диаметры эллипса, $O$ -- его центр. Проведем через точку $B$ перпендикуляр к прямой $OA$ и отложим на нем отрезки $BQ$ и $BQ$, равные $OA$. Доказать, что главные оси эллипса являются биссектрисами углов между прямыми $OP$ и $OQ$.
	\textit{Замечание.} Это дает способ построить главные оси эллипса по двум его сопряженным диаметрам.

	\item Пусть $a$ и $b$ -- фиксированные комплексные числа. Доказать, что при изменении $\varphi$ от $0$ до $2\pi$ точки вида $ae^{i\varphi} + be^{-i\varphi}$ заметают эллипс или отрезок.

	\item
	\begin{enumerate}
		\item На эллипсе взяты точки $A$ и $B$ так, что $\angle AOB  = \ang{90}$($O$ -- центр эллипса). Доказать, что $\dfrac{1}{OA^2} + \dfrac{1}{OB^2}$ не зависит от выбора точек $A$ и $B$.
		\item Доказать, что вписанные в эллипс ромбы описаны вокруго одной окружности.
	\end{enumerate}

	\item Объясните следующий парадокс. Эллипсы на плоскотси задают $5$-параметрическое семейство (эллипс задается двумя координатами центра $x_0$, $y_0$, двумя длинами полуосей $a$ и $b$ и углом $\varphi$ между осью $Ox$ и полуосью $a$). Наложим на параметры одно условие $a = b$. Почему в результате получается не $4$-параметрическое, а $3$-парамтерическое семейство (окружностей)?

	\item К эллипсу  с центром в $O$ проведены две параллельные касательные $l_1$ и $l_2$. Окружность с центром $O_1$ касается эллипса (внешним образом) эллипса и прямых $l_1$ и $l_2$. Доказать, что длина отрезка $OO_1$ равна сумме полуосей эллипса.
	\item Окружность радиуса $r$ с центром $C$, лежащим на большой полуоси эллипса касается эллипса(в двух точках); $O$ -- центр эллипса, $a$ и $b$ -- его полуоси. Доказать, что
	\[
		OC^2 = \frac{(a^2 - b^2)(b^2 - r^2)}{b^2}.
	\]

		\item Три окружности, центры которых лежат на большой оси эллипса, касаются эллипса. При этом окружность радиуса $r_2$ касается (внешним образом) окружностей радиусов $r_1$ и $r_3$. Доказать, что
	\[
		r_1 + r_3 = \frac{2a^2(a^2 - 2b^2)}{a^4} r_2.
	\]

	\item $N$ окружностей, центры которых лежат на большой оси эллипса, касаются эллипса. При этом окружность радиуса $r_i$ $(2 \leq i \leq N - 1)$ касается окружностей радиусов $r_{i-1}$ и $r_{i+1}$. Доказать, что если $3n - 2 \ge N$, то
	\[
		r_{2n-1}(r_n + r_{2n-1}) = r_n(r_n + r_{3n-2}).
	\]

	\item На комплексной плоскости дан треугольник $ABC$. Пусть $F_1$ и $F_2$ -- фокусы эллипса, касающегося сторон этого треугольника в их серединах. Доказать, что $F_1$ и $F_2$ -- корни производной многочлена третьей степени с корнями $A$, $B$ и $C$.
\end{enumerate}

\subsubsection*{4.1.5. Гипербола}

Решения задач $1$–$5$ аналогичны соответствующим задачам для эллипса.

\begin{enumerate}

	\item Доказать, что для гиперболы $\frac{x^2}{a^2} - \frac{y^2}{b^2} = 1$ существуют такие точка $F$ и прямая $d$, что отношение расстояний от точки гиперболы до точки $F$ и до прямой $d$ равно постоянному числу $e > 1$.

	\item Доказать, что разность расстояний от точки гиперболы до ее фокусов постоянно.

	\item Доказать, что в полярных координатах с начальным в фокусе гипербола имеет уравнение $r = \frac{p}{1 - e\cos\varphi}$, где $e > 1$.

	\item
	\begin{enumerate}
		\item Найти уравнение касательной к гиперболе $\dfrac{x^2}{a^2} - \dfrac{y^2}{b^2} = 1$ в точке $(x_0, y_0)$.
		\item Доказать, что касательная к гиперболе в точке $A$ является биссектрисой треугольника $F_1AF_2$.
	\end{enumerate}

	\item Доказать, что гипербола $\frac{x^2}{a^2} - \frac{y^2}{b^2} = 1$ является образом равнобочной гиперболы $x^2 - y^2 = 1$ при линейном преобразовании.

	\item Рассмотрим семейство гипербол $\frac{x^2}{a^2} - \frac{y^2}{b^2} = k$, зависящее от параметра $k$. Доказать, что середины всех хорд, которые эти гиперболы высекают на прямой $l$, совпадают.
	
	\item 
	\begin{enumerate}
		\item Касательная к гиперболе в точке $X$ пересекает асимптоты в точке $A$ и $B$. Доказать, что $X$ -- середина отрезка $AB$.
		\item Прямая перескает гиперболу в точках $A_1$ и $B_1$, а ее асимптоты -- в точках $A$ и $B$. Доказать, что $AA_1 = BB_1$.
	\end{enumerate}

	\item
	\begin{enumerate}
		\item Доказать, что площадь треугольника, образованного касательной к гиперболе и асимптотами, не зависит от выбора касательной.
		\item Рассмотрим параллелограмм, двумя противположными вершинами которого служат центр гиперболы и некоторая ее точка $X$. Доказать, что площадь параллелограмма не зависит от выбора точки $X$.
	\end{enumerate}

	\item Найдите множество точек пересечения пар перпендикулярных касательных к гиперболе.

	\item Вершины треугольника лежат на гиперболе $y = \frac{1}{x}$. Доказать, что его ортоцентр тоже лежит на этой гиперболе.

	\item Точка $M$ лежит на гиперболе $y = \frac{1}{x}$, а точка $M'$ симметрична ей относительно начала координат. Окружность радиуса $MM'$ с центром $M$ пересекает гиперболу помимо $M'$ еще и в точках $A$, $B$ и $C$. Доказать, что треугольник $ABC$ равносторонний.
\end{enumerate}

\section{4.2. Квадрики в $\mathbb{R}^n$. Проективная, аффинная и евклидова классификация квадрик}

\begin{enumerate}

	\item Доказать, что уравнение поверхности второго порядка, заданной в аффинном пространстве $\mathbb{R}^n$, переходом к новой системе координат (и, возможно, умножением обеих частей уравнения на $-1$) можно привести к одной из следующих нормальных форм:
	\begin{enumerate}
		\item $x_1^2 + \ldots + x_k^2 - x_{k+1}^2 - \ldots - x_r^2 = 1$, $0 \le k \le r \le n$;
		\item $x_1^2 + \ldots + x_k^2 - x_{k+1}^2 - \ldots - x_r^2 = 0$, $1 \le r \le n$, $k \ge \left[\dfrac{r+1}{2}\right]$;
		\item $x_1^2 + \ldots + x_k^2 - x_{k+1}^2 - \ldots - x_r^2 = 2x_{r + 1}$, $1 \le r \le n - 1$, $k \ge \left[\dfrac{r}{2}\right].$
	\end{enumerate}

	\item Найти максимальную размерность поверхности аффинного подпространства содержащихся в следующих поверхностях второго порядка в $R^n$:
	\begin{enumerate}
		\item $x_1^2 + \ldots + x_k^2 - x_{k+1}^2 - \ldots - x_n^2 = 0$;
		\item $x_1^2 + \ldots + x_k^2 - x_{k+1}^2 - \ldots - x_n^2 = 1$;
		\item $x_1^2 + \ldots + x_k^2 - x_{k+1}^2 - \ldots - x_{n-1}^2 = 2x_n$.
	\end{enumerate}

	\item Квадрика
	\[
		\sum_{i,j=1}^n a_{ij} x_i x_j + 2\sum_{i=1}^n b_i x_i + c = 0
	\]
	называется цилиндрической, если в некоторой системе координат ее можно задать уравнением с неполным числом переменных. Доказать, что квадрика цилиндрическая тогда и только тогда, когда
	\[
		\begin{vmatrix}
			a_{11} & \hdots & a_{1n} & b_1\\
			\hdots & \hdots & \hdots & \hdots\\
			a_{n1} & \hdots & a_{nn} & b_n \\
			b_1 & \hdots & b_n & c 
		\end{vmatrix} = 0,
		\quad
		\begin{vmatrix}
			a_{11} & \hdots & a_{1n} \\
			\hdots & \hdots & \hdots \\
			a_{n1} & \hdots & a_{nn}
		\end{vmatrix} = 0.
	\]

	\item
	\begin{enumerate}
		\item Доказать, что квадрика коническая тогда и только тогда, когда
		\[
			\mathrm{rank}\begin{pmatrix}
				a_{11} & \ldots & a_{1n} \\
				\vdots & \ddots & \vdots \\
				a_{n1} & \ldots & a_{nn}
			\end{pmatrix} =
			\mathrm{rank}\begin{pmatrix}
				a_{11} & \ldots & a_{1n} & b_1 \\
				\vdots & \ddots & \vdots & \vdots \\
				a_{n1} & \ldots & a_{nn} & b_n \\
				b_1 & \ldots & b_n & c
			\end{pmatrix}
		\]
		\item Доказать, что если для конической квадрики
		
		$\mathrm{rank}\begin{pmatrix}
			a_{11} & \ldots & a_{1n} \\
			\vdots & \ddots & \vdots \\
			a_{n1} & \ldots & a_{nn}
		\end{pmatrix} = n$, то она является конусом с единственной вершиной $0$. Если $M$ -- другая точка этого конуса, то ему принадлежат все точки прямой $OM$.
	\end{enumerate}

	\item Доказать, что квадрика является параболоидом тогда и только тогда, когда
	\[
		\begin{vmatrix}
			a_{11} & \hdots & a_{1n} \\
			\hdots & \hdots & \hdots \\
			a_{n1} & \hdots & a_{nn}
		\end{vmatrix} = 0, \quad
		\begin{vmatrix}
			a_{11} & \hdots & a_{1n} & b_1 \\
			\hdots & \hdots & \hdots & \hdots \\
			a_{n1} & \hdots & a_{nn} & b_n \\
			b_1 & \hdots & b_n & c
		\end{vmatrix} \ne 0.
	\]

	\item Доказать, что квадрику в евклидовом пространстве можно привести к одной из следующих форм:
	\begin{enumerate}
		\item $\dfrac{x_1^2}{a_1^2} + \ldots + \dfrac{x_k^2}{a_k^2} - \dfrac{x_{k+1}^2}{a_{k+1}^2} - \ldots - \dfrac{x_r^2}{a_r^2} = 1$;
		\item $\dfrac{x_1^2}{a_1^2} + \ldots + \dfrac{x_k^2}{a_k^2} - \dfrac{x_{k+1}^2}{a_{k+1}^2} - \ldots - \dfrac{x_r^2}{a_r^2} = 0$;
		\item $\dfrac{x_1^2}{p_1^2} + \ldots + \dfrac{x_k^2}{p_k^2} - \dfrac{x_{k+1}^2}{p_{k+1}^2} - \ldots - \dfrac{x_r^2}{p_r^2} = 2x_{r+1}$, $p_i > 0$.
	\end{enumerate}

	\item Для следующих евклидовых квадрик найдите канонические формы и канонические системы координат:
	\begin{enumerate}
		\item $2x_1 x_2 + 2x_1 x_3 - 2x_1 x_4 - 2x_2 x_3 - 2x_2 x_4 + 2x_3 x_4 - 2x_2 - 4x_3 - 6x_4 - 5 = 0$;
		\item $2x_1^2 + 3x_2^2 + 3x_3^2 + 3x_4^2 - 2x_1 x_2 - 2x_1 x_3 - 2x_1 x_4 - 2x_2 x_3 - 2x_3 x_4 - 2x_3 x_4 = 0$;
		\item $4x_1 x_2 + 4x_1 x_3 + 4x_1 x_4 + 4x_2 x_3 + 4x_2 x_4 + 4x_3 x_4 + 3x_4^2 + 14x_1 + 11 = 0$.
	\end{enumerate}

	\item Доказать, что вершины всех прямоугольных параллелепипедов, описанных вокруг эллипсоида, лежит на одной сфере.

\end{enumerate}

\section{4.3. Ответы, указания и решения}

\subsection{4.1. Кривые второго порядка}

\subsubsection**{4.1.1. Общие свойства кривых второго порядка (коник)}

1. \textit{Указание.} а) Пусть $a > b$. Тогда и для эллипса $\dfrac{x^2}{a-\lambda} + \dfrac{y^2}{b - \lambda} = 1$ и для гиперболы $\dfrac{x^2}{a-\lambda} - \dfrac{y^2}{\lambda - b} = 1$ фокус лежит на оси $Ox$ и удален от начала координат на расстояние
\[
	\sqrt{(a - \lambda) - (b - \lambda)} = \sqrt{(a - \lambda) + (\lambda - b)} = \sqrt{a - b}
\]
б) График функции $f(\lambda) = \dfrac{x^2}{a - \lambda} + \dfrac{y^2}{b - \lambda} = \sqrt{a - b}$ имеет вид
%Вставить картинку

Прямая $z = 1$ пересекает график ровно в двух точках, причем одна из них соотвествует эллипсу, а вторая -- гиперболе.

в) Если $F_1$ и $F_2$ -- фокусы эллипса и гиперболы, а $P$ -- их общая точка, то касательными к эллипсу и гиперболе будут биссектрисы внешнего и внутреннего углов $P$ треугольника $F_1PF_2$ треугольника соотвественно..

2. \textit{Указание.} Эллипс и гиперболу можно задать уравнением $A(z^2 + \overline{z}^2) + Bz \overline{z} = 1$ ($2A > B $ -- эллипс, $2A < B$ -- гипербола), а параболу уравнением $z^2 + \overline{z}^2 + 2z \overline{z} + 2iz - 2i \overline{z} = 0$. Будем считать, что уравнение коники имеет вид
\begin{equation}
	A(z^2 + \overline{z}^2) + Bz \overline{z} + Cz + \overline{C} \overline{z} + D = 0.
\end{equation}
Если $u$ -- центр правильного треугольника, вписанного в конику, то уравнению $(1)$ удовлетворяют точки $z = u + v \varepsilon^k$, $k = 0, 1, 2$, где $\varepsilon$ -- кубический корень из 1. Сложив три таких уравнения, получим
\begin{equation}
	A(u^2 + \overline{u}^2) + B(u \overline{u} + v \overline{v}) + Cu + \overline{C} \overline{u} + D = 0
\end{equation}
(мы воспользовались тем, что $\varepsilon^0 + \varepsilon^1 + \varepsilon^2 = 0$). Вычтем уравнение $(1)$ из уравнения $(2)$ при $z = u + v$. В результате получим $\operatorname{Re}(Fv + A \overline{v}^2) = 0$, где $F = 2Au + B \overline{u} + C$. Заменив $v$ на $v\varepsilon$ получим $Fv + A\overline{v}^2 = 0$. Так как $v \ne 0$, то $|v|^2 = |2Au + B \overline{u} + C|^2 A^{-2}$. Подставив это выражение в уравнение $(2)$, получим уравнение требуемой коники. 
\textit{Замечания.} 1. $A = 0$ соотвествует окружности.
2. Вторая коника совпадает с первой $\Leftrightarrow$ $B = 0$, т.е. в случае равнобочной гиперболы.
\subsubsection*{4.1.2. Теорема о пучке коник, проходящих через 4 точки}
3. \textit{Указание.} Можно считать, что $AB$ и $AD$ — оси координат $Ox$ и $Oy$, т.е. прямые $AB$ и $AD$ заданы уравнениями $y = 0$ и $x = 0$. Коэффициенты $\lambda$ и $\mu$ можно подобрать, рассмотрев ограничения левой и правой части на оси координат(ограничения на оси являются квадратными трехчленами с двумя общими корнями, поэтому они прпорциональны). Такой выбор $\lambda$ и $\mu$ означает, что
\[
	F(x, y) - \lambda y CD(x, y) - xBC(x, y) = xyQ,
\]
где $Q$ -- константа. В точке $C$ левя часть обращается в нуль, а $xy \ne 0$, поэтому $Q = 0$.

б) Согласно задаче $a)$ коники, проходящие черех 4 точки, образуют проективную прямую, порожденную точками $AB \cdot CD = 0$ и $AD \cdot BC = 0$; эта прямая порождена также парой любых других точек, например $f = 0 $ и $g = 0$.

4. \textit{Указание.} На направления осей коники влияют лишь квадратичные члены уравнения коники, поэтому будем учитывать только их. Можно считать, что уравнение одной из коник имеет вид $ax^2 + by^2 + \ldots = 0$ (для параболы $ab = 0$). Требуется выяснить, в каком случае линейная комбинация уравнений $ax^2 + by^2 + \ldots = 0$ и $a_1x^2 + b_1y^2 + c_1xy \ldots = 0$ может иметь вид $x^2 + y^2 + \ldots = 0$. Ясно, что $c_1 = 0$ — это необходимое условие. Оно означает, что оси коник перпендикулярны. Если $c_1 = 0$, то нужно подобратьт $\lambda$ так, чтобы выполнялось равенство $a + \lambda a_1 = b + \lambda b_1$, т.е. $\lambda = \frac{a - b}{b_1 - a_1}$($b_1 \ne a_1$ по условию). Плохой случай: $a + \lambda a_1 = b + \lambda b_1 = 0$, т.е. $a:b = a_1:b_1$. Рассматривая разность уравнений $ax^2 + by^2 + dx + ey + c =0$ и $ax^2 +by^2 + d_1 x + e_1 y + c_1 = 0$, легко проверить, что в этом случае коники имюет не более двух общих точек. 

5. \textit{Указание.}
 	а) квадратичная часть уравнения гиперболы с перпендикулярныим асимптотами имеет вид
	\[
		(x\cos\varphi + y\sin\varphi)(-x\sin\varphi + y\cos\varphi) = (y^2 - x^2)\sin\varphi \cos\varphi + xy(\cos^2 \varphi - \sin^2 \varphi).
	\]

	б) Согласно задаче а) линейная комбинация уравнений гиперболы с перпендикулярными асимптотами тоже является гиперболой с перпендикулярными асимптотами. В пучке коник , проходящих через точки $A$, $B$, $C$ и $H$, есть две (вырожденных) гиперболы с перпендикулярными асимптотами, а именно $AB \cdot CH = 0$ и $BC \cdot AH = 0.$ Следовательно, все коники этого пучка, являются гиперболами с перпендикулярными асимптотами.
 
6. \textit{Указание.} а) Через точки $K, L, M, N$ проходят три коники:
окружность $f = 0, g = KL \cdot MN = 0$ и $h = KN \cdot MN = 0$. Поэтому $h = \lambda f + \mu g$. Это равенство верно и для ограничений на прямую $AB$. Введем на прямой $AB$ координту $x$, причем $O$ -- начало координат. Тогда $f = x^2 - a, g = x^2$, поэтому $h = bx^2 - x$. Значти, корни уравнения $h = 0$ равноудалены от точки $O$.

б) решается аналогично задаче а).

7. \textit{Указание.} Коника $y^2 - \lambda x^2 + axy + cx + (b + \lambda)y + d = 0$ вырожденная тогда и только тогда, когда
$$
\begin{vmatrix}
	-\lambda & \frac{a}{2} & \frac{c}{2} \\
	\frac{a}{2} & 1 & \frac{b + \lambda}{2} \\
	\frac{c}{2} & \frac{b + \lambda}{2} & d 
\end{vmatrix} = 0.
$$
8. \textit{Указание.} а) и б) Пусть $f = 0$ -- уравнение данной коники. В неё вписаны четырёхугольники $ABCD$, $AFED$ и $BEFC$. Поэтому $f$ можно представить в виде:
\begin{align*}
	f &= \lambda_1 AB \cdot CD + \mu_1 AD \cdot BC, \tag{1} \\
	f &= \lambda_2 AF \cdot ED + \mu_2 AD \cdot FE, \tag{2} \\
	f &= \lambda_3 BE \cdot CF + \mu_3 BC \cdot EF. \tag{3}
\end{align*}
Приравнивая выражения $(1)$ и $(2)$, получаем
\[
	\lambda_1 AB \cdot CD - \lambda_2 AF \cdot ED = (\mu_1 BC - \mu_2 EF) AD.
\]
Левая часть обращается в нуль в точках пересечения прямых $AB$ и $ED$, $AF$ и $CD$. Следовательно, эти точки лежат на прямой $\mu_1 BC = \mu_2 EF$. На этой прямой лежит также точка пересечения прямых $BC$ и $EF$.
%поставить теги на все уравнения надо
Приравния $(3)$ и $(2)$, $(3)$ и $(1)$, получаем соответсвенно, что $\mu_2 AD = \mu_3 BC$ и $\mu_3 EF = \mu_1 AD$ -- прямые Паскаля шестиугольников $ADEBCF$ и $ADCFEB$. Легко проверить, что полученные три прямые пересекаются в одной точке, т.е. из любых двух уравнений следует третье (мы не рассматриваем вырожденные случаи, когда приходится сокращать на нуль).
в) Исходя из четырёхугольников $ABFE$, $ABDC$ и $CDEF$, представьте $f$ в трёх видах:
\begin{align*}
	f &= \lambda_1 AF \cdot BF + \mu_1 AB \cdot EF, \\
	f &= \lambda_2 AC \cdot BD + \mu_2 CD \cdot AB, \\
	f &= \lambda_3 CE \cdot DF + \mu_3 EF \cdot CD. 
\end{align*}
г) Переобозначение вершин шестиугольника не изменяет тройку прямых Штейнера, но изменяет тройку прямых Киркмана.
\subsubsection*{4.1.3 Парабола}
1. \textit{Указание.} Касательная имеет уравнение $y = 2x_0 x - x_0^2(y_0 = x_0^2)$.

2. \textit{Указание.} Прямая $y = ax + b$ пересекает параболу $y = x^2$ в точках $(x_1, y_1)$ и $(x_2, y_2)$, где $ax_i + b = x_i^2$; сумма корней этого квадратного уравнения равна $a$.

3. \textit{Указание.} Пусть прямая $l$ проходит через середины двух параллельных хорд параболы. Ось параллельна $l$ и проходит через середину хорды параболы, перпендикулярной $l$.

4. а) \textit{Указание.} Если $x_0^2 = 4ay_0$, то $x_0^2 + (y_0 - a)^2 = (y_0 + a)^2$.
б) \textit{Указание.} Пусть $(x - x_0)^2 + (y - y_0)^2 = (px + qy + r)^2$ для всех точек параболы $x^2 = 4ay$, причем $p^2 + q^2 = 1$. Выразив $y$ через $x$ мы должны получить тождество. Поэтому $b = 1, a = 0, x_0 = 0$ю Теперь легко проверить, что $v_0 = a, r = a.$
5. \textit{Указание.} Пусть касательная в точке $A(x_0, 4ax_0^2)$, пересекает ось $Oy$ в точке $B(0, -4ax_0^2)$. Тогда $AH = 4ax_0^2 + a = BF$, т.е. $BFAH$ - параллелограмм. Кроме того $FA = AH$.

6. \textit{Указание.} Расстояние от точки параболы с полярными координатами $(r, \varphi)$ до директрисы равно $r\cos\varphi + 2a$, а расстояние от точки параболы $B(0, -4ax_0^2)$ до параболы равно $r\cos\varphi + 2a$, поэтому $r = r\cos\varphi + 2a$.

7. \textit{Укзаание.} Точка $\cos\varphi + i\sin\varphi$ переходит в точку $\frac{-\cos\varphi + i\sin\varphi}{1 - \cos\varphi}$.

8. \textit{Ответ.} Существует.

\textit{Указание.} Положите $x = (x_1 + y_1)^2 + x_1y = x_1 + y_1$. Проверьте, что $x_1$ и $y_1$ можно выразить через $x$ и $y$.

 и $y$ можно выразить через $x'$ и $y'$.


9. \textit{Указание.} Подставив $y = x^2$ в уравнение $(x - x_0)^2 + (y - y_0)^2 = R^2$, получим уравнение 4-й степени с нулевым коэффициентом при $x^3$. Сумма корней этого уравнения равна нулю.


10. \textit{Указание.} Сложите уравнения $x^2 + py = 0$ и $(y - y_0)^2 + ax + b = 0$.

11. \textit{Указание.} Точки $P_1$ и $P_2$ имеют полярные координаты $(r+1, \varphi)$ и $(r_2, \varphi + \pi)$, где $r_1 = \frac{2a}{1 - \cos \varphi}$, $r_2 = \frac{2a}{1 - \cos(\varphi + \pi)} = \frac{2a}{1 + \cos \varphi}$ (задача 6).

12. \textit{Указание.} Пусть $A_1$ и $B_1$ — проекции точек $A$ и $B$ на директрису. Тогда $AO$ и $BO$ — серединные перпендикуляры к отрезкам $A_1F$ и $B_1F$. Поэтому $A_1O = FO = B_1O$, а значит, $\angle A_1B_1O = \angle B_1A_1O$. Следовательно, $\angle OFA = \angle OFB = 90^\circ - \varphi$, откуда $\angle AFB = 180^\circ - 2\varphi = \angle A_1 O B_1 = 2\angle AOB$.

13. \textit{Указание.} а) следует из задачи 12. б) Пусть $A_1$ и $B_1$ — проекции точек $A$ и $B$ на директрису. Точка $O$ — центр описанной окружности треугольника $A_1FB_1$. Она лежит на прямой $A_1B_1$ тогда и только тогда, когда $\angle AOB = 90^\circ$.

14. \textit{Указание.} Касательная в точке $(2t_i, t_i^2)$ задается уравнением $y = xt_i - t_i^2$.

15. \textit{Указание.} а) Пусть $A', B', C'$ — проекции фокуса $F$ на прямые $BC, CA, AB$. Точки $A', B', C'$ лежат на одной прямой (касательной, перпендикулярной оси параболы). Это означает, что точка $F$ лежит на описанной окружности треугольника $ABC$. В самом деле,
$$\angle AF C' = \angle AB'C' = \angle A'B'C = \angle A'FC\text{, поэтому}$$
$$\angle CFA = \angle A'FC = 180^\circ - \angle B.$$

б) Касательные к параболе $x^2 = 4y$ в точках $(2t_i, t_i^2)$ задаются уравнениями $y = t_ix - t_i^2$. Они пересекаются в точках $t_1 + t_2, t_1t_2)$. Легко проверить, что точка $(t_1+t_2+t_3 + t_1 t_2 t_3, -1)$ является ортоцентром треугольника с вершинами в этих точках.

в) Для параболы $x^2 = 4y$ и точек $\alpha, \beta, \gamma = (2t_i, t_i^2)$ получим:
\[S_{\alpha \beta \gamma} =
\begin{vmatrix} 
	2t_1 & t_1^2 & 1 \\
	2t_2 & t_2^2 & 1 \\
	2t_3 & t_3^2 & 1 
\end{vmatrix}, \quad S_{ABC} =
\frac{1}{2} \begin{vmatrix}
	t_2 + t_3 & t_2 t_3 & 1 \\
	t_1 + t_3 & t_1 t_3 & 1 \\
	t_1 + t_2 & t_1 t_2 & 1 
\end{vmatrix}
\]

16. \textit{Указание.} Касательные к параболе $x^2 = 4y$ в точках $(2t_1, t_1^2)$, $(2t_2, t_2^2)$ пересекаются в точке $(t_1+t_2, t_1 t_2)$. В нашем случае $t_1t_2 = -2$. Теперь легко проверить, что вершина параболы — ортоцентр треугольника $OAB$.

\subsubsection*{4.1.4. Эллипс}

1. \textit{Указание.} При $q = pe^2$ уравнение $\frac{(x - q)^2 + y^2}{(x - q)^2} = e^2$ эквивалентно уравнению $\frac{x^2}{p^2e^2} + \frac{y^2}{p^2 e^2(1 - e^2)} = 1$. При этом $1 - e^2 = \frac{b^2}{a^2}$, $p = \pm \frac{a}{e}$, $q = pe^2$; точка $F$ имеет координаты $(q, 0)$, а прямая $d$ задается уравнением $x = p$.

2. \textit{Указание.} Пусть $d_1$ и $d_2$ — директрисы, $H_1$ и $H_2$ — проекции точки $X$ на директрисы. Тогда $F_1X + F_2X = eH_1X + eH_2X = e(H_1X + H_2X) = e H_1 H_2$.

3. \textit{Указание.} Если $p$ — расстояние от фокуса до директрисы, то $\frac{r}{r \cos\varphi + a} = e$.

4. а) \textit{Ответ:} $\frac{xx_0}{a^2} + \frac{y y_0}{b^2} = 1.$

б) \textit{Указание.} Достаточно доказать, что если $l$ - биссектриса внешнего угла $A$ треугольника $F_1AF_2$, то $l$ пересекает эллипс лишь в точке $A$. Предположим, что $M$ — другая точка прямой $l$, а точка $F'_2$ симметрична $F_2$ относительно $l$. Тогда $F_1M + F_2M = F1_1M + F'_2M > F_1F_2' = F_1A + F_2A$, т.е. точка $M$ лежит вне эллипса.

5. \textit{Указание.} Рассмотрите преобразование $x \to ax, y \to by$.

6. \textit{Указание.} Представьте эллипс как образ окружности при линейном преобразовании.

7. \textit{Указание.} Любой параллелограмм является образом квадрата при аффинном преобразовании, а любой треугольник -- образом правильного треугольника.

8. \textit{Указание.} а) Эллипс является ортогональной проекцией окружности, а при такой проекции отношение площадей фигур сохраняется.

б) Можно считать, что точки $O, A$ и $B$ имеют координаты $(0, 0)$, $(a \cos\varphi, b\sin\varphi)$ и $(-a\sin\varphi, b\cos\varphi)$.


9. \textit{Указание.} Пусть $M$ — центр эллипса, $P_1$ и $P_2$ — проекции фокусов на касательную, $A$ — точка касания. Тогда согласно задаче 4,б $\angle P_1 A F_1 = \angle P_2 A F_2 = \phi$. Положим $x = F_1 A, y = F_2 A$. Тогда
\[
	P_1M^2 = P_2M^2 = (\frac{x+y}{2} \cos\varphi)^2 + (\frac{x+y}{2}\sin\varphi)^2 = \text{const}.
\]
Ясно также, что
\[
	d_1 d_2 = xy \sin^2 \phi
\]
и
\[
	F_1 F_2^2 = x^2 + y^2 + 2xy \cos 2\phi = (x + y)^2 - 4xy \sin^2 \phi,
\]
т.е. величина $xy \sin^2 \phi$ постоянна. 

Второе решение задачи б). Пусть $A$ — точка касания, точка $B$ симметрична $A$ относительно центра эллипса, точки $G_1$ и $H_1$ симметричны $F_1$ относительно касательных в точках $A$ и $B$. Тогда точки $G, H$ лежат на окружности радиуса $F_1 A + F_2 A = \text{const}$ с центром $F_2$. Поэтому величина $G_1 F_1 \cdot F_1 H_1 = 4d_1 d_2$ постоянна.

\textit{Третье решение задачи б).} Пусть касательные в точках $A$ и $A'$ пересекаются в точке $O$. Воспользовавшись задачей 10, получим
\[
	d_1 = OF_1 \sin \varphi, d_2 = OF_2 \sin(\alpha - \varphi), d_1' = OF_1 \sin(\alpha - \varphi), d_2' = OF_2 \sin \varphi.
\]

10. \textit{Указание.} Пусть точки $G_1$ и $G_2$ симметричны $F_1$ и $F_2$ относительно прямых $OA$ и $OB$ соответственно. Точки $F_1, B$ и $G_2$ лежат на одной прямой и $F_1 G_2 = F_1 B + B G_2 = F_1 B + B F_2$. Теугольники $G_2 F_1 O$ и $G_1 F_2 O$ имеют равные стороны. Следовательно, $\angle G_1 O F_1 = \angle G_2 O F_2$ и $\angle A F_1 O = \angle A G_1 O = \angle B F_1 O$. 

11. \textit{Указание.} Пусть $O$ — вершина рассматриваемого прямоугольника, $OA$ и $OB$ — касательные к эллипсу. Рассматриваемый в задаче 10 треугольник $F_1 O G_2$ в данном случае прямоугольный. Следовательно, $F_1 O^2 + F_2 O^2 = F_1 G_2^2 = \text{const}$. Если $M$ — центр эллипса, то $OM^2 = \frac{1}{4}(2 F_1 O^2 + 2 F_2 O^2 - F_1 F_2^2)$ — постоянная величина.

12. \textit{Указание.} Точка $Q$ лежит на описанной окружности треугольника $A F F_1$; при этом $Q$ — середина дуги $F_1 F_2$. Если $R$ — радиус описанной окружности, а $\alpha$ и $\beta$ — углы при вершинах $F_1$ и $F_2$, то $AQ = 2R \cos^2 \frac{\alpha - \beta}{2}, AP = 2R \sin^2 \frac{\alpha - \beta}{2}$. Поэтому $AP \cdot AQ = \left(2R \sin \frac{\alpha + \beta}{2} \cos \frac{\alpha + \beta}{2}\right)^2 = (R \sin \alpha + R \cos \beta)^2 = \left(\frac{AF_1 + AF_2}{2}\right)^2$.

13. \textit{Указание.} Точки $A$ и $B$ имеют координаты $(a \cos \varphi, b \sin \varphi)$ и $(a \sin \varphi, -b \cos \varphi)$. Точки $P$ и $Q$ имеют координаты $((a + b) \sin \varphi, -(a + b) \cos \varphi)$ и $((a - b) \sin \varphi, (a - b) \cos \varphi)$.

14. \textit{Указание.} Рассмотрим отображение $f(z) = az + b \overline z$. В координатах $(x, y)$, где $z = x + iy$, это отображение линейное, причем определитель равен $|a|^2 - |b|^2$. Образом окружности $|z| = 1$ при невырожденном линейном отображении будет эллипс, а при вырожденном -- отрезок или точка.

15. \textit{Указание.} а) Для эллипса $\frac{x^2}{a^2} + \frac{y^2}{b^2} = 1$ прямые $OA$ и $OB$ имеют уравнения $y = kx$ и $y = -\frac{x}{k}$. Точки $A$ и $B$ имеют координаты $(x_0, y_0)$ и $(x_1, y_1)$, где $x_0^2\left(\frac{1}{a^2} + \frac{k^2}{a^2}\right) = 1$ и $x_1^2\left(\frac{1}{a^2} + \frac{1}{k^2b^2}\right) = 1$. Поэтому 
\[
	\frac{1}{OA^2} + \frac{1}{OB^2} = \frac{\frac{1}{a^2} + \frac{k^2}{a^2}}{1 + k^2} + \frac{\frac{1}{a} + \frac{1}{k^2 + b^2}}{1 + \frac{1}{k^2}}= \frac{1}{a^2} + \frac{1}{b^2}.
\]
б) Если параллелограмм $ABCD$ вписан в эллипс, то прямые, соединяющие середины хорд $AB$ и $CD$, $BC$ и $AD$, проходят через центр эллипса. Следовательно, центр вписанного в эллипс ромба совпадает с центром эллипса. Радиус $R$ вписанной окружности ромба $ABCD$ равен высоте прямоугольного треугольника $AOB$, т.е. $\frac{1}{R^2} = \frac{1}{OA^2} + \frac{1}{OB^2}$.

16. \textit{Указание.} Самое поверхностное объяснение таково: при данном ограничении два из трех параметров $a, b$ и $\varphi$ становятся лишними. Чтобы получить более детальное объяснение, запишем уравнение эллипса, соответствующее параметрам $x_0, y_0, a, b, \varphi$:
\[
\frac{(x \cos \varphi + y \cos \varphi - x_0)^2}{a^2} + \frac{(-x_0 \sin \varphi + y \cos \varphi - y_0)^2}{b^2} = 1, 
\]
т.е. $px^2 + 2qxy + ry^2 + \dots = 0$, где $p = b^2 \cos^2 \varphi + a^2 \sin^2 \varphi$, $r = b^2 \sin^2 \varphi + a^2 \cos^2 \varphi$, $q = 2\cos\varphi\sin\varphi (a^2 - b^2) = (a^2 - b^2)\cos2\varphi$. Легко проверить, что $p - r = (b^2 - a^2) \cos 2\varphi$, поэтому $(p-r)^2 + q^2 = (a^2-b^2)^2$. Таким образом, из условия $a=b$ следуют сразу два условия: $p=r$ и $q=0$.

17. \textit{Указание.} Пусть $\varphi$ — угол между касательными $l_i$ и осью $Ox$. Рассмотрим окружность $S$ с центром $((a+b)\cos\varphi, (a+b)\sin\varphi)$, проходящую через точку $A = (a\cos\varphi, b\sin\varphi)$. Касательная к эллипсу в точке $A$ имеет уравнение $\frac{ \cos\varphi}{a^2}x + \frac{\sin\varphi}{b^2}y = 1$; прямая $AO_1$ имеет уравнение $y = \frac{a \sin\varphi}{b \cos\varphi} x + c$. Поэтому окружность $S$ касается эллипса. 

Если прямая $l_1$ касается эллипса в точке $(-a \cos\alpha, b \sin\alpha)$, то она имеет уравнение $-\frac{x \cos\alpha}{a} + \frac{y \sin\alpha}{b} = 1$, а значит, $\tg \varphi = \frac{b \cos\alpha}{a \sin\alpha}$. Квадрат расстояния от начала координат до прямой $l_1$ равен 
\[
	\frac{b^2}{sin^2 \alpha} \cdot \frac{1}{1 + \tg^2 \varphi} = b^2 \frac{\cos^2 \varphi}{\sin^2 \alpha} =
\]
\[
	= b^2 \cos^2 \varphi \left(1 + \frac{a^2 \sin^2 \varphi}{b^2cos^2 \varphi}\right) = b^2 \cos^2\varphi + a^2 \sin^2\varphi.
\]
Таким образом, окружность $S$ касается и прямых $l_1$ и $l_2$.

18. \textit{Указание.} Нормаль к эллипсу в точке $(x_0, y_0)$ имеет уравнение $\frac{-y_0}{b^2}x + \frac{x_0}{a^2}y = x_0 y_0 \left(\frac{1}{a^2} - \frac{1}{b^2}\right)$. Она пересекает большую полуось в точке с координатой $x_1 = b^2 x_0 \left(\frac{1}{b^2} - \frac{1}{a^2}\right) = x_0 \frac{a^2 - b^2}{a^2}$. При этом
\[
	r^2 = y_0^2 + x_0^2 \frac{b^4}{a^4} = b^2 - b^2 x_0^2\left(\frac{1}{a^2} - \frac{b^2}{a^4}\right).	
\]
Легко проверить, что $\frac{(a^2 - b^2)(b^2 - r)}{b^2} = x_i^2$.

19. \textit{Указание.} Соотношение
\[
	r_1 + r_2 = OC_1 - OC_2 = \frac{\sqrt{a^2 - b^2}}{b}(\sqrt{b^2 - r^2}- \sqrt{b^2 -r_2^2})
\]
можно преобразовать к виду:
\[
	a^4 r_1^2 - 2a^2(a^2 - 2b^2)r_1 r_2 + a^4 r_2^2 - 4b^4(a^2 - b^2) = 0. 
\]
Рассмотрим полученное выражение как уравнение относительно $r_1$. Оно имеет корни $r_1$ и $r_2$, поэтому $r_1 + r_2 = \frac{a^2(a^2 - 2b^2)}{a^4} r_2$.

20. \textit{Указание.} Согласно задаче $19$ числа $r_i$ удовлетворяют рекуррентному соотношению $r_{i+2} - k r_{i+1} + r_i = 0$, поэтому $r_p = a \lambda_1^p + b \lambda_2^p$, где $\lambda_1$ и $\lambda_2$ -- корни уравнения $x^2 - kx + 1 = 0$. При этом $\lambda_1 \lambda_2 = 1$, значит, $r_p = a \lambda^p + b \lambda^{-p}$. Требуемая формула проверяется очевидным образом.

21. \textit{Указание.} Можно считать, что треугольник $ABC$ получен преобразованием
\[
	z \to \frac{z + \overline z}{2} + \frac{z - \overline z}{2} \cos \alpha = z \cos^2 \frac{\alpha}{2} + \overline z \sin^2 \frac{\alpha^2}{2}
\]
из правильного треугольника с вершинами $w, \varepsilon w$ и $\varepsilon^2 w$, где $|w| = 1$ и $\varepsilon = \exp(2\pi i / 3)$. Тогда полуоси $a$ и $b$ рассматриваемого эллипса равны $\frac{1}{2}$ и $\frac{\cos\alpha}{2}$, а расстояние между его фокусами $F_1$ и $F_2$ равно $\sqrt{a^2 - b^2} = \frac{1}{2} \sin \alpha$. 
Точки $F_1$ и $F_2$ переходят в точки $(\pm 1, 0)$ при растяжении с коэффициентом $(\sin\alpha)^{-1} = \left(\sin \frac{\alpha}{2} \cos \frac{\alpha}{2}\right)^{-1}$ . В результате такого растяжения получаем преобразование $z \to z \ctg \alpha + \overline z \tg \alpha$. Положим $a = w \ctg \alpha$. Тогда многочлен с корнями $A, B$ и $C$ имеет вид
\[
	P(x) = (x - a - \frac{1}{a})(x - a \varepsilon - \frac{1}{a \varepsilon})(x - a \varepsilon^2 - \frac{1}{a \varepsilon^2}).
\]
Легко проверить, что $P'(x) = 3x^2 + 3\varepsilon + 3 \overline \varepsilon = 3x^2 - 3$, поэтому корни многочлена $P'$ равны $\pm 1$.

\subsubsection*{4.1.5. Гипербола}

4. а) \textit{Ответ:} $\dfrac{xx_0 }{a^2} - \dfrac{yy_0}{b^2} = 1$.

б) \textit{Указание.} Если $B$ — другая точка биссектрисы угла $F_1AF_2$, то $|F_1B - F_2B| < F_1 F_2' = |F_1A - F_2A|$.

6. \textit{Указание.} Точки пересечения прямой $y = px + q$ и гиперболы $\frac{x^2}{a^2} - \frac{y^2}{b^2} = k$ можно найти, решив уравнение.
\[
	\frac{x^2}{a^2} - \frac{(px + q)^2}{b^2} = \left(\frac{b^2 - a^2 p^2}{b^2}\right)x^2 - \frac{2px}{b^2} + q^2 = k.
\]
Полусумма корней этого уравнения не зависит от $k$. Тем самым, координата $x$ середины хорды постоянна.

7. \textit{Указание.} Решим сначала задачу б). Асимптоты гиперболы $\frac{x^2}{a^2} - \frac{y^2}{b^2} = 1$ имеют уравнение $\frac{x^2}{a^2} - \frac{y^2}{b^2} = 0$. Поэтому согласно задаче $6$ середины отрезков $AB$ и $A_1B_1$ совпадают. В том случае, когда $A_1 = B_1$, получаем задачу а).

8. \textit{Указание.} а) Касательная $\frac{xx_0}{a^2} - \frac{yy_0}{b^2} = 1$ пересекает асимптоты $y = \pm \frac{b}{a}x$ в точках с координатами $x_{1,2} = a\left(\frac{x_0}{a} \mp \frac{y_0}{b}\right)^{-1}$, поэтому $x_1 x_2 = a^2$. 

б) Любую гиперболу можно получить линейным преобразованием из гиперболы $y = x^{-1}$, для которой утверждение задачи очевидно. 

\textit{Замечание.} Согласно задаче 7, а) отношение площадей рассматриваемых треугольника и параллелограмма равно $2$.

9. \textit{Ответ.} Пусть $\varphi$ — величина того угла между асимптотами, который содержит гиперболу. При $\varphi \ge 90^\circ$ ГМТ пусто, при $\varphi < 90^\circ$ — окружность с центром в центре гиперболы, из которой выброшены 4 точки пересечения с асимптотами. 

\textit{Указание.} Воспользуйтесь решением задачи 11 темы "Эллипс".

10. \textit{Указание.} Пусть $a = \alpha + \frac{i}{\alpha}, b = \beta + \frac{i}{\beta}$ и $c = \gamma + \frac{i}{\gamma}$ — вершины данного треугольника (на комплексной плоскости). Проверьте, что $h = - \frac{1}{\alpha \beta \gamma} - i \alpha \beta \gamma$ — его ортоцентр. Покажите, например, что число $\frac{a - h}{b - c}$ чисто мнимое.

11. \textit{Указание.} Пусть $M = (a, \frac{1}{a})$. Координаты $x$ точек $M', A, B$ и $C$ удовлетворяют уравнению $x^4 - 2ax^3 + \dots = 0$. Следовательно, $x_A + x_B + x_C +(-a) = 2a$, т.е. $x_A + x_B + x_C = 3x_M$. Аналогично $y_A + y_B + y_C = 3y_M$. Т.е. $\vec{OA} + \vec{OB} + \vec{OC} = 3\vec{OM}$ и $\vec{MA} + \vec{MB} + \vec{MC} = \vec{0}$.

\end{document}
